\begin{figure*}[!htbp]
\begin{paperbox}{ScienceWorld 实例提示词}
\texttt{<task\_description>} \\
\texttt{\{\{task\_description\}\}} \\
\texttt{</task\_description>}

\vspace{0.2cm}
\texttt{<initial\_observation>} \\
\texttt{\{\{initial\_observation\}\}} \\
\texttt{</initial\_observation>}

\vspace{0.2cm}
\texttt{<instructions>} \\
完成上文描述的科学实验。你正在与模拟环境交互。请一次发出一条命令并观察结果。

\vspace{0.2cm}
\textbf{提示}：使用 \texttt{?navigation} 或 \texttt{?object} 等查询命令探索可用动作，且不消耗轮次。

\vspace{0.2cm}
\textbf{策略提示：}
\begin{itemize}
\item 先探索：使用 \texttt{open door to [room]}，再使用 \texttt{go to [room]} 导航
\item 观察每个房间，寻找所需对象
\item 使用 \texttt{pick up [object]} 拾取对象
\item 加热时：找到炉灶，用 \texttt{activate [stove]} 打开，将容器放到其上
\item 冷却时：使用冷冻柜或冰箱
\item 使用 \texttt{focus on [object]} 检查物质
\end{itemize}

\vspace{0.2cm}
为完成任务，请执行任务描述所需的动作。当你认为任务完成时，发出命令：

\vspace{0.2cm}
\texttt{```text} \\
\texttt{task completed} \\
\texttt{```} \\

\vspace{0.2cm}
\textbf{请记住：}
\begin{itemize}
\item 逐步思考需要哪些动作
\item 探索环境以找到所需对象
\item 有些动作可能需要前提条件（例如先拾取对象才能使用它）
\end{itemize}
\texttt{</instructions>}
\end{paperbox}
\end{figure*}
